\subsection{Von Neumann rank}
We can construct an hierarchy of sets, well ordered with respect with $\subseteq$, 
called the Von Neumann hierarchy, in such a way that any well-formed set in ZF is subset of 
at least one of the sets in this hierarchy. We can define the Von Neumann rank of a set $A$ in 
terms of the Von Neumann hierarchy as follows:

\begin{definition}[Von Neumann universe]
    For every ordinal $i$ we define the Von Neumann universe $V_i$ as the set which 
    obey the following construction:
    
    \begin{equation*}
        V_i = \begin{cases}
            \varnothing            & \text{if } i = 0 \\
            \bigcup P(V_{j < i})   & otherwise \\
        \end{cases}
    \end{equation*}
\end{definition}

\begin{theorem}[Von Neumann theorem]
    Every set $S$ is subset of some universe $V_i$
\end{theorem}
\begin{proof}
    Every entity in set theory is a set, moreover. Every set is either empty or is a set containing 
    other sets. Let's proceed by induction. The base case is
    $V_0$ which is the empty set. Assume all the element of some set $S$ are subsets of  
    some Von Neumann universe $V_j$, we can assert that $S$ is subset of $V_{j+1}$ by 
    the definition of $V_{j+1}$.
\end{proof}

\begin{definition}[Von Neumann rank]
    For every set $S$ we can define $r(S)$ the Von Neumann rank of $S$ as the smallest
    ordinal $i$ such that $S \subseteq V_i$.
\end{definition}
